\bookchapter{Conclusion : Rester curieux, rester lucide}

Félicitations d'être allé jusqu'au bout.

Ce livre a démonté le matériel du smartphone jusqu'à l'écosystème logiciel, des réglages quotidiens au mode développeur, de la sécurité au patrimoine numérique. S'il vous a aidé à résoudre un seul problème concret — ne serait-ce que désactiver une autorisation superflue ou prendre une fois l'habitude de sauvegarder — ces mots n'ont pas été écrits en vain.

Mais plus important que toute technique, deux choses :

\begin{itemize}
  \item \textbf{Rester curieux}: Le smartphone évolue chaque année, l'IA évolue, les pliables évoluent, la communication par satellite évolue. Les étapes apprises aujourd'hui seront peut-être obsolètes demain, mais la curiosité de « vouloir savoir pourquoi » ne sera jamais dépassée. Poser une question de plus — « à quoi sert cette autorisation », « où sont stockées ces données », « cette fonction va-t-elle fuiter ma vie privée » — est plus utile que de mémoriser cent procédures.
  \item \textbf{Rester lucide}: Le smartphone est un outil, pas un maître. Il est conçu pour capter votre attention, occuper votre temps et guider votre consommation. L'utiliser est votre choix ; être utilisé par lui est une perte de contrôle. Regardez régulièrement votre temps d'écran, demandez-vous « suis-je en train d'utiliser le smartphone pour faire quelque chose, ou le smartphone me consume-t-il ? ». Ce n'est pas de la sensiblerie, c'est un instinct de survie à l'ère numérique.
\end{itemize}

La technologie change, mais les besoins humains ne changent pas — nous avons besoin de sécurité, de connexion et de respect. Le smartphone n'est qu'un support. Le véritable « cours sur le smartphone » n'est pas dans ce livre, mais à l'instant où vous déverrouillez l'écran chaque jour.

\begin{tcbpractice}{Dernier exercice}
Fermez ce livre, prenez votre smartphone, et faites une chose :\textbf{Prenez 30 secondes pour réfléchir : quelles sont les trois choses dans votre smartphone dont la perte vous ferait vraiment mal ?}Puis assurez-vous qu'elles sont sauvegardées. Ce livre aura alors valu la peine.

\end{tcbpractice}
